\section{Proofs for the concrete formulations}
\label{app:formulations}

\subsection{Grouped paraboloids and arbitrary PSD column packing}
For a grouped tangent $(\dot s,\dot w)$, put
$\dot q=\dot s-2\ip{w}{\dot w}$. Direct differentiation gives
\begin{equation}
 D^2(-\log q)[(\dot s,\dot w),(\dot s,\dot w)]
       =(\dot q/q)^2+2\norm{\dot w}^2/q.
 \label{eq:grouped-metric}
\end{equation}
For a packed block, an arbitrary affine tangent gives
$\dot D=\dot S-\dot W^TW-W^T\dot W$, and
\begin{equation}
 D^2(-\log\det D)[(\dot S,\dot W),(\dot S,\dot W)]
 =\tr((D^{-1}\dot D)^2)
       +2\tr(D^{-1}\dot W^T\dot W).
 \label{eq:packing-metric}
\end{equation}
These are Hessians in affine $(s,w)$ or $(S,W)$ coordinates; treating
$D$ as an independent affine variable when taking the second derivative
would omit the second term. With $c_\ell$ columns in block $\ell$,
trace Cauchy--Schwarz gives
\begin{equation}
 \norm{\dot z}_F^2\geq
 \sum_\ell\frac1{c_\ell}
            \left(\frac{\mathrm d}{\mathrm dt}\log\det D_\ell\right)^2.
 \label{eq:packing-contraction}
\end{equation}
Its scalar version is the simultaneous contraction onto every $\log q$.
Both Hessians are positive definite on the free affine coordinates.
Standard self-concordance follows by restricting the classical PSD
barrier to the affine matrices
$\left(\begin{smallmatrix}I&W\\W^T&S\end{smallmatrix}\right)$, whose
determinant is $\det D$; the one-column case gives \eqref{eq:grouped-metric}.
The derivative $DF=-\sum_\ell\tr(D_\ell^{-1}\dot D_\ell)$ and
\eqref{eq:packing-contraction} imply
$|DF[\dot z]|\leq\sqrt H\norm{\dot z}_F$, where
$H=\sum_\ell c_\ell=kh$. This proves the parameter upper bound $H$
without charging the fixed identity blocks.

For a feasible target write $\delta_a=1-\ip{c_a}{w_a}$.
The diagonal residuals $d_\gamma=(D_{\ell(\gamma)})_{\gamma\gamma}$
satisfy
\[
 \sum_{\gamma\text{ from }a}d_\gamma=1-\norm{w_a}^2
 \leq1-\ip{c_a}{w_a}^2\leq2\delta_a.
\]
The same identity holds for grouped $q_{aG}$. If the target error is
at most $\eps$, then $\sum_a\delta_a\leq\eps$ and
\begin{equation}
 \prod_\ell\det D_\ell
 \leq\prod_\gamma d_\gamma
 \leq\prod_a(2\delta_a/h)^h
 \leq(2\eps/H)^H.
 \label{eq:packing-collapse}
\end{equation}
The first inequality is Hadamard's inequality and permits every
completion variable to vary; it is not a diagonal-fiber assumption.
Integrating \eqref{eq:packing-contraction} and using weighted
Cauchy--Schwarz shows that every joining curve has length at least
$H^{-1/2}$ times the absolute change in the sum of log determinants.
If the reference has $D_\ell=q_0I$,
\eqref{eq:packing-collapse} gives \eqref{eq:grouped-packed-distance}.

For completeness, these are indeed the claimed centers. At fixed $W$
and fixed residual diagonals, Hadamard's inequality minimizes the
barrier precisely when every residual $D_\ell$ is diagonal.
Partial minimization over the unrestricted off-diagonal entries of $S$
therefore leaves the grouped problem. Its equations in $s_{aG}$ give
one common residual $q$ per ball, and its equations in $w$ give
$w_{aG}=(\tau q/2)c_{aG}$. The source equation becomes
$hq+(\tau q/2)^2=1$, proving \eqref{eq:grouped-packed-center}.
Positive definite Hessians make these stationary points the unique
centers. Differentiated central stationarity on the affine slice yields
$H_\tau z'(\tau)=\nabla\ell$ and hence
\[
 \left\|\frac{dz}{d\log\tau}\right\|^2
 =\tau^2\frac{\mathrm d}{\mathrm d\tau}\ell(z(\tau))
 =k\tau^2r'(\tau)
 =H\left(1-\frac h{\sqrt{h^2+\tau^2}}\right).
\]
At a center, $DF=\tau\,d\ell$ on the slice, so the same expression is
the squared dual gradient norm. Its limit $H$ proves the exact parameter.
Changing variables from $\tau$ to $r$ gives the length element
$\sqrt H\rho'(r)\dd r$, proving \eqref{eq:formulation-route}.

\subsection{Linked norm-tree centers and distance}
The positive-sheet domain \eqref{eq:norm-tree-domain} is a convex affine
restriction of a product of Lorentz cones. Its Hessian is positive
definite: any nonzero free direction changes at least one cone block,
whose Hessian is positive definite. Restriction and addition preserve
standard self-concordance. Telescoping cancels every nonroot internal
square and gives, for each source,
\begin{equation}
 \sum_vq_{av}=1-\norm{w_a}^2.
 \label{eq:tree-telescoping}
\end{equation}
The derivative in a nonroot axial variable is
$-2t_{av}/q_{av}+2t_{av}/q_{a,\mathrm{parent}(v)}$.
Since $t_{av}>0$, central stationarity forces all residuals in a tree
to agree. Leaf stationarity gives $w_a=rc_a$, and
\eqref{eq:tree-telescoping} becomes $bq+r^2=1$, with $r=\tau q/2$.
The positive internal coordinates are uniquely
$t_{av}=(n_vq+r^2C_{av})^{1/2}$; their root value is one.
These equations prove feasibility and stationarity, and strict convexity
proves uniqueness. They give the analytic center $w=0,t_{av}^2=n_v/b$
and the path/speed formulas asserted in the main text.

For each Lorentz block, the parameter-two gradient inequality gives
$(d\log q_{av})^2\leq2D^2(-\log q_{av})$ after restriction.
Therefore the total metric controls the Euclidean log-residual metric
with factor $1/\sqrt2$. More sharply, with the exact restricted
parameter $\nu$ proved below, the barrier itself is
$\sqrt\nu$-Lipschitz in metric length. Telescoping and the same source
AM--GM calculation as \eqref{eq:packing-collapse} give
\[
 \prod_{a,v}q_{av}(y)\leq(2\eps/L)^L,\qquad
 F(y)-F(\bar z)\geq L\log(k/(2\eps)).
\]
This proves \eqref{eq:tree-distance}. Differentiated stationarity gives
the same speed as for grouped balls with $h=b,H=L$, so the same scalar
integration proves \eqref{eq:formulation-route}. Thus no asymptotic
regularity of boundary KKT equations is needed for the matched route.

\subsection{Exact root-slice parameter}
We prove \eqref{eq:tree-parameter} for one tree. Before fixing the root
$t$, let $\Phi=-\sum_v\log q_v$ on the linked cone and $H=\Phi''$.
Its logarithmic degree is $2b$; hence
$H z=-\Phi'(z)$ and $\norm{\Phi'}_{H^{-1}}^2=2b$.
Let $e_t$ be the root-coordinate covector. Orthogonal projection of the
gradient onto $\{h:e_t^Th=0\}$ gives the exact restricted norm
\begin{equation}
 \norm{DF}_{F,*}^2=2b-\frac{t^2}{e_t^TH^{-1}e_t},\qquad
 \nu_T=2b-\frac1{\sup_z\beta_T(z)},\qquad
 \beta_T(z)=\frac{e_t^TH^{-1}e_t}{t^2}.
 \label{eq:root-leverage}
\end{equation}
Dikin containment and the halfspace $t>0$ imply
$e_t^TH^{-1}e_t\leq t^2$, so $\beta_T\leq1$.
Indeed, if the inequality failed, a vector of local norm less than one
could move the root coordinate to zero.

If a root child is a free leaf, set $t=1$ and that leaf to
$\sqrt{1-\delta}$. Set other root leaves to zero and scale every
internal child subtree from a fixed interior point by $\delta^2$.
The root residual is $q=\delta-O(\delta^4)>0$.
For the ambient trial direction $h_t=1$, $h_{\rm leaf}=1$, and every
other component zero, descendant blocks contribute zero, while
\[
 D^2\Phi[h,h]
 =\frac{\{2(1-\sqrt{1-\delta})\}^2}{q^2}\longrightarrow1.
\]
Here the root residual's second derivative in that direction is zero.
The variational identity
$(e_t^TH^{-1}e_t)^{-1}=\min_{h_t=1}h^THh$ now gives
$\sup\beta_T=1$. This includes a one-node tree without descendants.
It proves $\nu_T=2b-1$ with an explicit feasible limiting sequence.

Suppose instead that all root children are internal, with positive
axial coordinates $a_1,\ldots,a_c$. Eliminate all their strict
descendants from the Hessian. The effective axial curvature $h_i$ of
child subtree $i$ is at least $a_i^{-2}$ by the same Dikin halfspace
argument on that subtree. Increasing an axial curvature increases the
root Schur complement, so a lower bound is obtained from
\[
 -\log(t^2-\textstyle\sum_i a_i^2)-\sum_i\log a_i.
\]
Normalize $t=1$, put $r=\sum_i a_i^2$, $q=1-r$, and
$Z=\sum_i a_i^4/(q+2a_i^2)$. Its child Hessian is a diagonal matrix
with entries $2/q+a_i^{-2}$ plus $4aa^T/q^2$; the root cross row is
$-4a^T/q^2$, and the root entry is $2(1+r)/q^2$.
The rank-one inverse formula therefore gives the root Schur complement
\begin{equation}
 S=\frac{2(1+r)-16Z/(q+4Z)}{q^2}.
 \label{eq:root-schur}
\end{equation}
Since $x/(q+2x)$ increases in $x\geq0$,
$\sum_i a_i^4/(q+2a_i^2)\leq r^2/(q+2r)=r^2/(1+r)$.
The inequality $S\geq2$ is equivalent to
$r(3-r)\geq4(2-r)Z$; replacing $Z$ by this upper bound leaves
$3r(1-r)^2/(1+r)\geq0$. Thus $\beta_T\leq1/2$.
Conversely, scale all child subtrees by $\delta\downarrow0$.
The trial direction $h_t=1$, with all other variables fixed, has norm
squared $2(1+r)/(1-r)^2\to2$.
Together with the lower bound on the minimizing Schur complement this
proves $\sup\beta_T=1/2$ and $\nu_T=2b-2$.
For independent trees, gradient norm squares and their suprema add,
proving the product parameter used in \eqref{eq:tree-distance}.

\subsection{Heterogeneous source counts and objective weights}
A single weighted version records the dependence on source sizes for
all three formulations. Let source $a$ have $m_a$ groups or internal
nodes, let $M_0=\sum_am_a$, $p_a=m_a/M_0$, and maximize
$\sum_a\lambda_a\ip{c_a}{w_a}$ with every $\lambda_a>0$.
Define the scale
\begin{equation}
 \Delta_{\rm eff}=\prod_a(\lambda_a/p_a)^{p_a}.
 \label{eq:formulation-entropy}
\end{equation}
At the separate analytic centers all source residuals are $1/m_a$.
For an accurate target, put
$e_a=\lambda_a(1-\ip{c_a}{w_a})$, so $\sum_a e_a\leq\eps$.
Telescoping, or the residual diagonal identities and Hadamard's
inequality, show that the ratio of terminal to central residual products
is at most
$\prod_a(2e_a/\lambda_a)^{m_a}$.
This product is maximized at $e_a=\eps p_a$, giving
\[
 F(y)-F(\bar z)\geq M_0\log(\Delta_{\rm eff}/(2\eps)).
\]
Consequently the entire accurate set is at distance at least
\begin{equation}
 \frac{M_0}{\sqrt\nu}\pos{\log(\Delta_{\rm eff}/(2\eps))},
 \quad
 \nu=\begin{cases}M_0,&\text{grouped or packed},\\
                   \sum_a\nu_{T_a},&\text{norm trees}.
       \end{cases}
 \label{eq:heterogeneous-formulation-distance}
\end{equation}
The grouped/packed parameter remains exact by applying the center
calculation source by source, with parameter $\tau\lambda_a$.
For equal weights $\lambda_a=\lambda>0$,
$\Delta_{\rm eff}=\lambda\exp(-\sum_ap_a\log p_a)$;
weighted AM--GM gives $\Delta_{\rm eff}\leq\sum_a\lambda_a$ with
equality precisely when $\lambda_a$ is proportional to $m_a$.
Thus the scale, as well as the square-root count, matters for movement.
